\subsection{Classes of Simple Modules}

Denote by \(\mathrm{Is}_{\mathrm{A}}\{\mathrm{X},\mathrm{Y}\}\) the relation

``A is a ring and X, Y are isomorphic A-modules.''

This is an equivalence relation with respect to X and Y. If X is an A-module, then we denote the class of objects equivalent to X for \(\operatorname{Is}_{\mathrm{A}}\) by \(\operatorname{cl}(\mathrm{X})\) and call it the class of the A-module X (Set Theory, II, §6, No.~9, p.~122). By definition, \(\operatorname{cl}(\mathrm{X})\) is an A-module isomorphic to X; moreover, two A-modules X and Y are isomorphic if and only if we have \(\operatorname{cl}(\mathrm{X}) = \operatorname{cl}(\mathrm{Y})\) .

Let A be a ring. The relation

``\(\lambda\) is a class of finitely generated A-modules''

is collectivizing in \(\lambda\) (Set Theory, II, §1, No.~4, p.~68). Indeed, every finitely generated A-module is isomorphic to an A-module of the form \(A_{s}^{n}/R\) , where n is a natural number and R a submodule of \(A_{s}^{n}\) , so that our assertion follows from Set Theory, II, §6, No.~9, p.~122.

We denote the set of classes of finitely generated A-modules by \(\mathcal{F}(A)\) . Every simple A-module is monogenous (VIII, p.~46, Proposition 1), and consequently the classes of the simple A-modules form a subset of \(\mathcal{F}(A)\) , which we from now on denote by \(\mathcal{S}(A)\) (or simply S). When the ring A is commutative, the mapping \(\mathfrak{m} \mapsto \operatorname{cl}(A/\mathfrak{m})\) is a bijection from the set of maximal ideals of A to the set \(\mathcal{S}(A)\) (loc. cit. and VIII, p.~46, Remark 1). When A is left Artinian, the set \(\mathcal{S}(A)\) is finite (VIII, p.~51, Remark b)).

